\subsection{Dijkstra con heapq (Caminos Mínimos desde un Origen)}

\graybold{Preguntas guía:}

\begin{itemize}
\item ¿Me piden la distancia mínima desde un nodo origen a todos los demás en un grafo con pesos NO negativos?
\item ¿Los pesos son variables, de modo que un BFS simple (que asume peso 1 por arista) no sirve?
\item ¿El grafo es disperso ($m$ del orden de $n$), con $n$ y $m$ de hasta unos pocos cientos de miles?
\item ¿Puede haber aristas paralelas, lazos, o varias formas de llegar a un mismo nodo que mejoran su distancia más de una vez?
\end{itemize}

\graybold{Utilidad:} Calcula las distancias mínimas desde un origen a todos los nodos cuando los pesos son no negativos, procesando los nodos en orden creciente de distancia con una cola de prioridad. Es el algoritmo base de caminos mínimos: rutas en mapas, costo mínimo en grafos de estados, y como subrutina de problemas más grandes (varios orígenes con un supernodo, grafo con estados extra como "¿ya usé el cupón?", reconstrucción del camino guardando el padre de cada nodo). Con pesos negativos deja de ser correcto y hace falta Bellman-Ford.

\graybold{Complejidad:} \bigO{(n + m) \log m} con la versión de borrado perezoso: cada arista puede insertar a lo sumo una entrada en el heap, así que el heap tiene \bigO{m} elementos en total.

\graybold{Fundamento teórico:} El invariante es que, cuando un nodo sale del heap con la menor distancia tentativa de todo el heap, esa distancia es definitiva. La razón es que los pesos son no negativos: cualquier otro camino hacia ese nodo tiene que pasar primero por algún nodo que todavía está en el heap, con distancia tentativa mayor o igual, y agregarle aristas de peso $\ge 0$ no puede hacerlo más corto. Por eso los nodos se fijan en orden creciente de distancia, y al fijar $u$ basta relajar sus aristas salientes: si $\text{dist}[u] + w < \text{dist}[v]$, se mejora $\text{dist}[v]$ y se inserta $(\text{dist}[v], v)$ en el heap. Como \texttt{heapq} no permite disminuir la prioridad de un elemento ya insertado, se usa BORRADO PEREZOSO: se inserta una entrada nueva con la distancia mejorada y la vieja queda en el heap; cuando una entrada sale con distancia mayor que $\text{dist}[u]$, es obsoleta y se descarta con un \texttt{continue}. Así un nodo puede estar varias veces en el heap, pero solo se procesa una vez, y el total de inserciones está acotado por el número de relajaciones exitosas, es decir por $m$. La hipótesis de pesos no negativos es imprescindible: con una arista negativa, un nodo ya fijado podría mejorar después, y el invariante se rompe.~\cite{dijkstra1959}

\graybold{Ejercicio aplicado}

(CSES 1671 — Shortest Routes I) Dado un grafo dirigido de $n$ ciudades y $m$ vuelos con longitud, calcular la distancia mínima desde la ciudad 1 a cada ciudad (se garantiza que todas son alcanzables).

\graybold{I/O:} $n \le 10^5$, $m \le 2\cdot10^5$ con tres enteros por arista $\Rightarrow$ lectura total + tokenización (patrón 2). Las distancias pueden llegar a \aprox{}\ten{14}, sin problema para los enteros de Python. La salida son $n$ números en una sola línea, unidos con un único \texttt{write}. Verificado contra el caso de ejemplo y contrastado con Bellman-Ford en 400 tandas aleatorias con lazos, aristas paralelas y pesos de 1 a \ten{9}, sin discrepancias. Con $n = 10^5$ y $m = 2\cdot10^5$ tarda entre \aprox{}0.24s y \aprox{}0.35s en los tres grafos medidos: árbol más aristas aleatorias, una cadena de peso 1 con atajos caros desde el origen (cada nodo se inserta dos veces), y un grafo con muchas relajaciones sucesivas; el límite es de 1 segundo.

\graybold{Código solución}

\begin{lstlisting}[language=Python]
import sys, heapq

def solve():
    data = sys.stdin.buffer.read().split()
    n, m = int(data[0]), int(data[1])
    adj = [[] for _ in range(n + 1)]
    for i in range(m):
        a, b, c = int(data[2+3*i]), int(data[3+3*i]), int(data[4+3*i])
        adj[a].append((b, c))
    INF = float('inf')
    dist = [INF]*(n + 1)
    dist[1] = 0
    pq = [(0, 1)]
    while pq:
        d, u = heapq.heappop(pq)
        if d > dist[u]:
            continue          # entrada obsoleta: u ya salio con una distancia menor
        for v, w in adj[u]:
            nd = d + w
            if nd < dist[v]:
                dist[v] = nd
                # no se puede "bajar la prioridad" en heapq: se inserta de nuevo
                heapq.heappush(pq, (nd, v))
    sys.stdout.write(" ".join(map(str, dist[1:])) + "\n")

solve()
\end{lstlisting}

\graybold{Extras: versión con entradas codificadas en un entero} — aproximadamente un 20\% más rápida, útil si el grafo es más grande o el límite más ajustado. En lugar de tuplas, cada arista se guarda como \texttt{peso << 17 | destino} y cada entrada del heap como \texttt{distancia << 17 | nodo}: comparar enteros es más barato que comparar tuplas, y como la distancia ocupa los bits altos, el orden del heap sigue siendo por distancia. Requiere que $2^{\text{SH}}$ supere al mayor número de nodo; con $n \le 10^5$ basta $\text{SH} = 17$. Verificada igual que la versión principal; en los mismos grafos tarda entre \aprox{}0.22s y \aprox{}0.28s.

\begin{lstlisting}[language=Python]
import sys
from heapq import heappush, heappop

def solve():
    data = sys.stdin.buffer.read().split()
    n = int(data[0]); m = int(data[1])
    A = list(map(int, data[2:2 + 3*m:3]))
    B = list(map(int, data[3:3 + 3*m:3]))
    C = list(map(int, data[4:4 + 3*m:3]))
    SH = 17                          # 2^SH debe ser mayor que n
    MASK = (1 << SH) - 1
    adj = [[] for _ in range(n + 1)]
    for a, e in zip(A, [(c << SH) | b for b, c in zip(B, C)]):
        adj[a].append(e)             # arista: peso << SH | destino
    INF = 1 << 62
    dist = [INF]*(n + 1)
    dist[1] = 0
    pq = [1]                         # distancia 0, nodo 1  ->  (0 << SH) | 1
    while pq:
        x = heappop(pq)
        u = x & MASK                 # bits bajos: nodo
        d = x >> SH                  # bits altos: distancia
        if d > dist[u]:
            continue
        for e in adj[u]:
            v = e & MASK
            nd = d + (e >> SH)
            if nd < dist[v]:
                dist[v] = nd
                heappush(pq, (nd << SH) | v)
    sys.stdout.write(" ".join(map(str, dist[1:])) + "\n")

solve()
\end{lstlisting}
